\documentclass[12pt,a4paper]{article}
\usepackage[utf8]{inputenc}
\usepackage[T1]{fontenc}
\usepackage[brazil]{babel}
\usepackage{geometry}
\usepackage{listings}
\usepackage{xcolor}
\usepackage{hyperref}
\usepackage{titlesec}
\usepackage{parskip}
\usepackage{fancyhdr}
\usepackage{graphicx}
\usepackage{booktabs}
\usepackage{array}
\usepackage{mdframed}
\usepackage{amssymb}

\geometry{margin=2.5cm}

% Colors
\definecolor{codebg}{rgb}{0.95,0.95,0.97}
\definecolor{codeframe}{rgb}{0.7,0.7,0.8}
\definecolor{keyword}{rgb}{0.13,0.29,0.53}
\definecolor{string}{rgb}{0.16,0.52,0.16}
\definecolor{comment}{rgb}{0.5,0.5,0.5}
\definecolor{histbg}{rgb}{0.98,0.96,0.88}
\definecolor{histframe}{rgb}{0.85,0.75,0.4}
\definecolor{titlebg}{rgb}{0.13,0.29,0.53}

% Code style
\lstdefinestyle{kotlin}{
  backgroundcolor=\color{codebg},
  frame=single,
  rulecolor=\color{codeframe},
  basicstyle=\ttfamily\small,
  keywordstyle=\color{keyword}\bfseries,
  stringstyle=\color{string},
  commentstyle=\color{comment}\itshape,
  breaklines=true,
  breakatwhitespace=true,
  tabsize=4,
  showstringspaces=false,
  numbers=left,
  numberstyle=\tiny\color{comment},
  xleftmargin=1.5em,
  framexleftmargin=1.5em,
}
\lstset{style=kotlin}

% History box
\newmdenv[
  backgroundcolor=histbg,
  linecolor=histframe,
  linewidth=1.5pt,
  leftmargin=0,
  rightmargin=0,
  innerleftmargin=12pt,
  innerrightmargin=12pt,
  innertopmargin=8pt,
  innerbottommargin=8pt,
]{historybox}

% Header/Footer
\pagestyle{fancy}
\fancyhf{}
\rhead{\textcolor{comment}{\small GuardApp — Documentação Técnica}}
\lhead{\textcolor{comment}{\small Incognia Technical Challenge}}
\rfoot{\thepage}

% Title formatting
\titleformat{\section}{\Large\bfseries\color{titlebg}}{}{0em}{}[\titlerule]
\titleformat{\subsection}{\large\bfseries\color{keyword}}{}{0em}{}
\titleformat{\subsubsection}{\normalsize\bfseries}{}{0em}{}

\hypersetup{
  colorlinks=true,
  linkcolor=keyword,
  urlcolor=keyword,
}

\begin{document}

% ─── CAPA ────────────────────────────────────────────────────────────────────
\begin{titlepage}
  \centering
  \vspace*{2cm}
  {\Huge\bfseries\color{titlebg} GuardApp}\\[0.5cm]
  {\Large Android Environment Detection SDK}\\[0.3cm]
  {\large\color{comment} Documentação Técnica com Contexto Histórico}\\[2cm]
  \rule{\textwidth}{1pt}\\[0.5cm]
  {\normalsize
    \textbf{Projeto:} Incognia Technical Challenge\\[0.3cm]
    \textbf{Propósito:} Detecção de ambientes adulterados em Android\\[0.3cm]
    \textbf{Linguagem:} Kotlin / Android SDK\\[0.3cm]
    \textbf{Cobertura:} Emulador, Root, Clone App, Hooking, Reempacotamento
  }\\[2cm]
  \rule{\textwidth}{1pt}\\[1cm]
  {\small\color{comment}
    Cada classe, função e decisão de design é explicada junto\\
    com a história que motivou sua existência.
  }
  \vfill
  {\small 2024}
\end{titlepage}

\tableofcontents
\newpage

% ─── VISÃO GERAL ─────────────────────────────────────────────────────────────
\section{Visão Geral da Arquitetura}

O GuardApp é um SDK Android que detecta ambientes adulterados em tempo de execução.
Responde a uma pergunta crítica para aplicações financeiras e de segurança:

\begin{center}
\textit{``Este app está rodando num dispositivo legítimo, ou em ambiente controlado por um atacante?''}
\end{center}

\subsection{Estrutura de Pacotes}

\begin{lstlisting}
com.incognia.guardapp
  detection/
    EnvironmentAnalyzer.kt    # contrato (interface)
    DetectionSignal.kt        # tipos de dados
    SecurityReport.kt         # resultado final
    EnvironmentGuard.kt       # orquestrador (facade)
    DetectionSignatures.kt    # catalogo de assinaturas
    analyzers/
      EmulatorAnalyzer.kt     # detecta AVD/QEMU/Genymotion
      CloneAnalyzer.kt        # detecta VirtualApp/Parallel Space
      VirtualizationAnalyzer.kt # detecta Frida/Xposed
      RootDetector.kt         # detecta root/su/Magisk
      SignatureAnalyzer.kt    # verifica assinatura do APK
  ui/
    MainActivity.kt           # interface do usuario
\end{lstlisting}

\subsection{Fluxo de Execução}

\begin{lstlisting}
MainActivity
    |
    v
EnvironmentGuard.generateReport()
    |
    +-- async --> EmulatorAnalyzer.analyze()     ---+
    +-- async --> CloneAnalyzer.analyze()        ---+
    +-- async --> VirtualizationAnalyzer.analyze()--+--> flatten()
    +-- async --> RootDetector.analyze()         ---+        |
    +-- async --> SignatureAnalyzer.analyze()    ---+        v
                                                       SecurityReport
                                                            |
                                                    isTampered / riskScore
                                                    emulatorProbability
                                                    summary
\end{lstlisting}

Os quatro analisadores rodam em paralelo via coroutines Kotlin.
Qualquer falha individual é capturada silenciosamente --- o relatório
continua com os sinais dos demais analisadores.

\newpage

% ─── ARQUIVO 1 ───────────────────────────────────────────────────────────────
\section{EnvironmentAnalyzer.kt --- O Contrato}

\begin{lstlisting}
fun interface EnvironmentAnalyzer {
    fun analyze(context: Context): List<DetectionSignal>
}
\end{lstlisting}

\subsection{O que e}

Uma \texttt{fun interface} --- um contrato com um unico metodo.
Qualquer classe que quiser ser um analisador de ambiente implementa esse contrato.

\subsection{Por que comecar aqui}

Toda a arquitetura do projeto nasce desse contrato.
Antes de falar de emulador, root ou Frida, a pergunta que o projeto responde e:
\textit{``qual e a forma padrao de perguntar ao ambiente se algo esta errado?''}

\begin{historybox}
\textbf{Historia: fun interface e SAM Conversion}

O Java introduziu interfaces funcionais no Java 8 (2014) com a anotacao
\texttt{@FunctionalInterface}. A ideia veio das linguagens funcionais --- Haskell,
ML, Scala --- onde funcoes sao valores de primeira classe.

O Kotlin evoluiu o conceito com \texttt{fun interface} (SAM --- Single Abstract Method):
qualquer interface com um unico metodo pode ser instanciada com uma lambda diretamente,
sem precisar de classe anonima.

\begin{lstlisting}
// Sem fun interface (verbose):
val fake = object : EnvironmentAnalyzer {
    override fun analyze(ctx: Context) = listOf(sinal)
}

// Com fun interface (SAM):
val fake = EnvironmentAnalyzer { listOf(sinal) }
\end{lstlisting}

Nos testes, isso e fundamental --- criar analisadores falsos fica trivial.
\end{historybox}

\subsection{Decisoes de Design}

\begin{itemize}
  \item Recebe \texttt{Context} --- passaporte do Android para acessar PackageManager,
        ApplicationInfo, etc.
  \item Retorna \texttt{List<DetectionSignal>} e nao \texttt{Boolean} --- um
        analisador pode encontrar varios problemas simultaneamente.
  \item E uma \texttt{interface} e nao \texttt{abstract class} --- permite composicao
        sem heranca rigida.
\end{itemize}

\newpage

% ─── ARQUIVO 2 ───────────────────────────────────────────────────────────────
\section{DetectionSignal.kt --- Os Dados}

\begin{lstlisting}
enum class SignalCategory(val displayName: String) {
    EMULATOR("Emulator"),
    CLONE_APP("Clone / Dual-Space App"),
    VIRTUALIZATION("Virtualization / Hooking"),
    ROOT("Root / Superuser"),
    SIGNATURE("Signature / Integrity"),
    SUSPICIOUS_PATH("Suspicious Path"),
    SYSTEM_PROPERTY("Suspicious System Property"),
}

enum class Severity { LOW, MEDIUM, HIGH }

data class DetectionSignal(
    val category: SignalCategory,
    val description: String,
    val severity: Severity,
    val explanation: String = "",
)
\end{lstlisting}

\subsection{Os tres tipos}

\textbf{SignalCategory} --- de onde veio o sinal (qual analisador, qual ameaca).

\textbf{Severity} --- o peso do sinal:

\begin{center}
\begin{tabular}{lll}
\toprule
Valor & Significado & Contribuicao ao riskScore \\
\midrule
LOW & Informacional, nao prova nada sozinho & 0 \\
MEDIUM & Suspeito, aumenta o risco & 1 \\
HIGH & Forte indicador de adulteracao & 3 \\
\bottomrule
\end{tabular}
\end{center}

\textbf{DetectionSignal} --- o sinal em si com quatro campos.

\begin{historybox}
\textbf{Historia: data class no Kotlin}

Em Java, criar uma classe de dados exigia escrever manualmente \texttt{equals()},
\texttt{hashCode()}, \texttt{toString()} e \texttt{copy()} --- dezenas de linhas
de codigo boilerplate propensas a bugs.

A JetBrains criou o Kotlin em 2011 (lancado em 2016) com \texttt{data class}:
uma unica palavra-chave gera todos esses metodos automaticamente.

\begin{lstlisting}
// Java: ~40 linhas
// Kotlin:
data class DetectionSignal(
    val category: SignalCategory,
    val description: String,
    val severity: Severity,
)
\end{lstlisting}

Importante para os testes: \texttt{assertEquals(sinal1, sinal2)} funciona
por valor (compara campos), nao por referencia (compara ponteiros).
\end{historybox}

\newpage

% ─── ARQUIVO 3 ───────────────────────────────────────────────────────────────
\section{SecurityReport.kt --- O Resultado}

\begin{lstlisting}
data class SecurityReport(val signals: List<DetectionSignal>) {

    val isTampered: Boolean
        get() = signals.any { it.severity != Severity.LOW }

    val riskScore: Int
        get() = signals.sumOf { signal ->
            when (signal.severity) {
                Severity.HIGH   -> 3
                Severity.MEDIUM -> 1
                Severity.LOW    -> 0
            }.toInt()
        }

    val emulatorProbability: Float
        get() {
            var p = 0.0
            emulatorSignals.forEach { signal ->
                p += (1.0 - p) * when (signal.severity) {
                    Severity.HIGH   -> 0.80
                    Severity.MEDIUM -> 0.50
                    Severity.LOW    -> 0.10
                }
            }
            return p.toFloat().coerceIn(0f, 1f)
        }
}
\end{lstlisting}

\subsection{isTampered}

Se existe qualquer sinal MEDIUM ou HIGH, o ambiente e considerado adulterado.
LOW sozinho e apenas informacional.

\subsection{riskScore}

HIGH vale 3 pontos, MEDIUM vale 1. Permite comparar gravidade entre dispositivos.
Um score 9 e muito mais grave que score 3, mesmo que ambos sejam \texttt{isTampered = true}.

\subsection{emulatorProbability --- A Matematica Bayesiana}

\begin{historybox}
\textbf{Historia: Thomas Bayes e Probabilidade Independente (1763)}

Thomas Bayes, matematico ingles, publicou em 1763 (postumamente) o teorema
que leva seu nome. A ideia central: \textbf{evidencias independentes se acumulam}.

A formula usada no codigo:
$$P = 1 - \prod_{i}(1 - p_i)$$

Na pratica: cada novo sinal contribui com a probabilidade restante de ainda
nao ter sido detectado:

\begin{center}
\begin{tabular}{lll}
\toprule
Sinal & Calculo & Resultado \\
\midrule
1 HIGH (0.80) & $0 + (1-0) \times 0.80$ & $0.80$ \\
2 HIGH (0.80) & $0.80 + (1-0.80) \times 0.80$ & $0.96$ \\
3 MEDIUM (0.50) & $0.96 + (1-0.96) \times 0.50$ & $0.98$ \\
\bottomrule
\end{tabular}
\end{center}

Por que nao somar? Porque $0.80 + 0.80 + 0.50 = 2.10$ e impossivel
(probabilidade nao pode ultrapassar 1.0).
\end{historybox}

\newpage

% ─── ARQUIVO 4 ───────────────────────────────────────────────────────────────
\section{EnvironmentGuard.kt --- O Orquestrador}

\begin{lstlisting}
class EnvironmentGuard private constructor(private val context: Context) {

    private val analyzers: List<EnvironmentAnalyzer> = listOf(
        EmulatorAnalyzer(),
        CloneAnalyzer(),
        VirtualizationAnalyzer(),
        SignatureAnalyzer(),
        RootDetector(),
    )

    suspend fun generateReport(): SecurityReport = coroutineScope {
        val deferred = analyzers.map { analyzer ->
            async(Dispatchers.Default) {
                runCatching { analyzer.analyze(context) }.getOrDefault(emptyList())
            }
        }
        SecurityReport(deferred.awaitAll().flatten())
    }

    companion object {
        fun create(context: Context): EnvironmentGuard =
            EnvironmentGuard(context.applicationContext)
    }
}
\end{lstlisting}

\subsection{private constructor + companion object}

O construtor e privado --- ninguem cria \texttt{EnvironmentGuard(context)} diretamente.
A unica forma e via \texttt{EnvironmentGuard.create(this)}.

Isso forca o uso de \texttt{context.applicationContext} em vez do
\texttt{Activity context}, evitando memory leak --- a Activity pode ser destruida,
o \texttt{applicationContext} nao.

\subsection{Execucao Paralela com Coroutines}

\begin{historybox}
\textbf{Historia: Coroutines e Concorrencia}

O conceito de coroutines foi proposto por \textbf{Melvin Conway} em 1963 --- o mesmo
Conway da Lei de Conway. A ideia: funcoes que podem suspender e resumir execucao
sem bloquear uma thread.

O Kotlin implementou coroutines como biblioteca de primeira classe em 2018.
\texttt{async/await} permite paralelismo sem callbacks ou thread pools manuais:

\begin{lstlisting}
// Sem paralelismo (sequencial - mais lento):
analyzers.forEach { it.analyze(context) }

// Com async/await (paralelo):
val deferred = analyzers.map { async(Dispatchers.Default) { it.analyze(context) } }
deferred.awaitAll()  // espera todos terminarem
\end{lstlisting}

Cada analisador roda em \texttt{Dispatchers.Default} --- um pool de threads
otimizado para operacoes de CPU. File I/O e network probes nao bloqueiam uns aos outros.
\end{historybox}

\subsection{runCatching --- Isolamento de Falhas}

\begin{lstlisting}
runCatching { analyzer.analyze(context) }.getOrDefault(emptyList())
\end{lstlisting}

Se um analisador crasha, o erro e capturado silenciosamente e retorna lista vazia.
Um analisador quebrado nao derruba o relatorio inteiro.

\newpage

% ─── ARQUIVO 5 ───────────────────────────────────────────────────────────────
\section{EmulatorAnalyzer.kt --- Detectando Emuladores}

\begin{lstlisting}
override fun analyze(context: Context): List<DetectionSignal> =
    checkBuildProperties() +
    checkSystemProperties() +
    checkEmulatorFiles() +
    checkCpuInfo()
\end{lstlisting}

\begin{historybox}
\textbf{Historia: QEMU e o AVD (2003-2007)}

\textbf{Fabrice Bellard}, engenheiro frances, criou o \textbf{QEMU} em 2003 ---
emulador de processador completo que executa instrucoes ARM num computador x86.
Bellard e famoso na comunidade por projetos extremos: tambem criou o FFmpeg.

O Google pegou o QEMU, colocou o Android por cima, e chamou de \textbf{AVD}
(Android Virtual Device) em 2007. Os engenheiros escolheram ``goldfish'' como
nome do perfil de hardware virtual. Ninguem pensou que isso seria um ponto
fraco de seguranca duas decadas depois.
\end{historybox}

\subsection{Camada 1: checkBuildProperties()}

\begin{lstlisting}
val hardware = Build.HARDWARE.lowercase()
if (hardware in setOf("goldfish", "ranchu", "vbox86", "nox"))
    signal("Emulator hardware detected: ${Build.HARDWARE}")
\end{lstlisting}

\texttt{Build.*} sao constantes gravadas em \texttt{/system/build.prop} pelo
fabricante durante a compilacao. Em emuladores, a Google nao se preocupa em
disfarcar:

\begin{center}
\begin{tabular}{ll}
\toprule
Campo & Emulador vs Device Real \\
\midrule
\texttt{Build.HARDWARE} & \texttt{goldfish} vs \texttt{shiba} (Pixel 8 Pro) \\
\texttt{Build.FINGERPRINT} & comeca com \texttt{generic} vs hash do fabricante \\
\texttt{Build.MODEL} & \texttt{Android SDK built for x86} vs \texttt{Pixel 8 Pro} \\
\texttt{Build.BRAND} & \texttt{generic} vs \texttt{google} \\
\bottomrule
\end{tabular}
\end{center}

\subsection{Camada 2: checkSystemProperties() via Reflection}

\begin{lstlisting}
val clazz = Class.forName("android.os.SystemProperties")
val get = clazz.getMethod("get", String::class.java, String::class.java)
val value = get.invoke(null, "ro.kernel.qemu", "")
\end{lstlisting}

\begin{historybox}
\textbf{Historia: Reflection no Java (1997)}

A Sun Microsystems adicionou Reflection ao Java 1.1 em 1997 como parte do
Java Core Reflection API. A ideia: inspecionar e invocar qualquer classe
e metodo em runtime, ignorando visibilidade (\texttt{private}, \texttt{internal}).

\texttt{android.os.SystemProperties} e uma API \textbf{interna} do Android ---
nao esta no SDK publico, nao compila se voce chamar diretamente.
Via Reflection, voce burla essa restricao em runtime.
\end{historybox}

\subsection{Camada 3: checkEmulatorFiles()}

\begin{lstlisting}
val EMULATOR_PATHS = listOf(
    "/dev/socket/qemud",       // socket QEMU host<->guest
    "/dev/qemu_pipe",          // pipe de alta performance
    "/dev/socket/genyd",       // Genymotion
    "/system/bin/nox-prop",    // Nox
)
EMULATOR_PATHS.filter { File(it).exists() }
\end{lstlisting}

Arquivos de dispositivo criados pelo kernel do emulador no boot.
Em hardware fisico, simplesmente nao existem.

\subsection{Camada 4: checkCpuInfo() com charArrayOf}

\begin{lstlisting}
// "goldfish" nunca aparece como literal no APK compilado:
val goldfish = charArrayOf('g','o','l','d','f','i','s','h').concatToString()
val cpuInfo = File("/proc/cpuinfo").readText()
if (cpuInfo.contains(goldfish)) { ... }
\end{lstlisting}

\begin{historybox}
\textbf{Historia: /proc/cpuinfo e a tecnica charArrayOf}

Tom Killian criou o \texttt{/proc} no Bell Labs em 1984. O arquivo
\texttt{/proc/cpuinfo} expoe o tipo de CPU do hardware diretamente do kernel.
Em emuladores QEMU, aparece ``QEMU Virtual CPU'' e ``goldfish'' ---
o kernel do emulador se autoidentifica.

A tecnica \texttt{charArrayOf} vem da pratica de obfuscacao de malware:
se voce escrever \texttt{"goldfish"} como string literal, ela aparece em
texto puro no APK e qualquer um executa \texttt{strings apk.apk | grep goldfish}.
Construindo em runtime a partir de chars individuais, a string nunca
existe como literal no binario.
\end{historybox}

\newpage

% ─── ARQUIVO 6 ───────────────────────────────────────────────────────────────
\section{CloneAnalyzer.kt --- Detectando Apps Clone}

\begin{lstlisting}
class CloneAnalyzer(
    private val uidProvider: () -> Int = { Process.myUid() },
    private val fileProbe: (String) -> Boolean = { File(it).exists() },
    private val packageChecker: (Context, String) -> Boolean = { ctx, pkg ->
        runCatching { ctx.packageManager.getPackageInfo(pkg, 0); true }
            .getOrDefault(false)
    },
    private val dataDirProvider: (Context) -> String? = { it.applicationInfo?.dataDir },
) : EnvironmentAnalyzer
\end{lstlisting}

\begin{historybox}
\textbf{Historia: VirtualApp e o mercado de clones (2015)}

Em 2015, um estudante universitario chines conhecido como \textbf{Lody} (GitHub: asLody)
publicou o \textbf{VirtualApp} --- open source, em Java. A pergunta que motivou:
\textit{``E se eu criasse um Android dentro do Android?''}

O VirtualApp intercepta chamadas Binder (o protocolo IPC do Android, criado por
\textbf{Dianne Hackborn} baseado no OpenBinder da Be Inc. em 1996) usando Reflection
para substituir os proxies locais por versoes proprias.

Em 2016, a \textbf{LBE Tech} lancou o \textbf{Parallel Space} baseado na mesma
tecnica --- 50 milhoes de downloads em meses. Bancos viram isso em fraudes
a partir de 2017: o clone framework interceptava camera, GPS e verificacoes
de seguranca antes do app financeiro acessar.
\end{historybox}

\subsection{Camada 3: UID Offset --- A Deteccao Mais Elegante}

\begin{lstlisting}
val userId = uidProvider() / 100_000
if (userId != 0) {
    signals += DetectionSignal(
        SignalCategory.CLONE_APP,
        "App running under user ID $userId (non-primary user)",
        Severity.HIGH
    )
}
\end{lstlisting}

\begin{historybox}
\textbf{Historia: Android Multi-User e UIDs (2012)}

Ken Thompson e Dennis Ritchie criaram o sistema de UIDs no Unix em 1969.
No Android 4.2 (2012), a Google adicionou suporte a multiplos usuarios
para tablets compartilhados por familias. A decisao de engenharia: alocar
UIDs em blocos de 100.000 por usuario.

\begin{center}
\begin{tabular}{ll}
\toprule
Usuario & Range de UIDs \\
\midrule
0 (primario) & 10.000 -- 99.999 \\
1 (secundario) & 110.000 -- 199.999 \\
2 & 210.000 -- 299.999 \\
\bottomrule
\end{tabular}
\end{center}

\textbf{O clone app precisa de usuario secundario.} Nao ha como escapar.
A decisao de engenharia de 2012 para familias em tablets virou,
12 anos depois, a deteccao mais confiavel contra clone apps.

\begin{lstlisting}
// Device limpo:    myUid() = 10.153  / 100.000 = 0  -> sem sinal
// Dentro do clone: myUid() = 110.153 / 100.000 = 1  -> HIGH
\end{lstlisting}
\end{historybox}

\subsection{Camada 5: /proc/self/maps e .so baked-in}

\begin{lstlisting}
File("/proc/self/maps").forEachLine { line ->
    if (line.contains("io.va", ignoreCase = true)) {
        signals += DetectionSignal(...)
    }
}
\end{lstlisting}

\begin{historybox}
\textbf{Historia: ELF, shared libraries e paths baked-in (1988)}

O formato \textbf{ELF} (Executable and Linkable Format) foi introduzido
no Unix System V em 1988. Shared libraries (\texttt{.so}) tem seus paths
de origem gravados pelo linker dentro do proprio binario compilado --- o \textbf{rpath}.

Um atacante pode renomear o APK do VirtualApp com \texttt{apktool}.
Mas o \texttt{libva.so} dentro do APK e um binario compilado --- o \texttt{apktool}
nao toca nele. O path \texttt{com.lody.virtual} continua baked-in no binario.

O \texttt{/proc/self/maps} --- criado por Tom Killian no Bell Labs em 1984 ---
expoe quais \texttt{.so} estao carregadas na memoria. O kernel nao mente.
\end{historybox}

\subsection{Injecao de Dependencias para Testabilidade}

\begin{historybox}
\textbf{Historia: Dependency Injection e TDD (Michael Feathers, 2004)}

\textbf{Michael Feathers} formalizou em 2004 no livro \textit{``Working Effectively
with Legacy Code''}: qualquer codigo que nao pode ser testado sem o ambiente
real e legacy code, independente de quando foi escrito.

A solucao: substituir chamadas ao sistema por parametros injetaveis.
No Kotlin moderno, uma lambda com valor default resolve sem frameworks:

\begin{lstlisting}
// Producao: usa sistema real (nao passa nada)
val analyzer = CloneAnalyzer()

// Teste: simula UID 110.000 (usuario secundario)
val analyzer = CloneAnalyzer(uidProvider = { 110_042 })
\end{lstlisting}
\end{historybox}

\newpage

% ─── ARQUIVO 7 ───────────────────────────────────────────────────────────────
\section{VirtualizationAnalyzer.kt --- Detectando Hooking}

\begin{lstlisting}
override fun analyze(context: Context): List<DetectionSignal> =
    checkHookFiles() +
    checkProcMaps() +
    checkXposedClassLoader() +
    checkFrida()
\end{lstlisting}

\subsection{checkHookFiles() --- Artefatos de Cada Framework}

\begin{historybox}
\textbf{Historia: Xposed Framework (rovo89, Alemanha, 2012)}

rovo89 era um desenvolvedor alemao frustrado com a impossibilidade de modificar
apps do sistema Android sem recompilar o Android inteiro. A solucao: substituir
\texttt{/system/bin/app\_process} --- o programa que inicia o Zygote e e pai
de todo app Android --- por uma versao que injetava \texttt{XposedBridge.jar}
no classpath de todo processo.

\medskip
\textbf{Historia: LSPosed (comunidade, 2020)}

O Xposed modificava \texttt{/system} diretamente, quebrando o \textbf{dm-verity}
(verificacao criptografica de integridade da Google). O LSPosed usou Magisk
como base --- monta um overlay virtual sem tocar no sistema real.

\medskip
\textbf{Historia: Cydia Substrate (Jay ``saurik'' Freeman, EUA, 2008)}

Jay Freeman criou o Cydia em 2008 --- loja alternativa para iPhone com jailbreak.
O Cydia Substrate usava \texttt{LD\_PRELOAD} (feature do linker Unix) para forcar
o carregamento de sua biblioteca antes de qualquer outra.

\medskip
\textbf{Historia: Magisk (topjohnwu / Wang Guang, Taiwan, 2016)}

Estudante taiwanes criou o Magisk: \textbf{systemless root} --- monta \texttt{tmpfs}
(filesystem em RAM) sobre \texttt{/system}. O Android ve modificacoes.
O sistema real esta intacto. O dm-verity fica feliz.
\end{historybox}

\subsection{checkXposedClassLoader() --- O Classpath como Detector}

\begin{lstlisting}
try {
    Class.forName("de.robv.android.xposed.XposedBridge")
    signals += DetectionSignal(..., Severity.HIGH)  // encontrou = Xposed ativo
} catch (_: ClassNotFoundException) { /* device limpo */ }
\end{lstlisting}

\begin{historybox}
\textbf{Historia: ClassLoader e Lazy Loading (James Gosling, 1995)}

James Gosling projetou o Java em 1995 com \textbf{lazy loading} --- classes
so existem quando pedidas. O ClassLoader busca no classpath quando
\texttt{Class.forName()} e chamado.

O Xposed injeta \texttt{XposedBridge.jar} no classpath do \textbf{Zygote}.
Todo processo filho herda o classpath do pai. Logo,
\texttt{de.robv.android.xposed.XposedBridge} fica acessivel em qualquer app.

Se \texttt{Class.forName()} nao lanca \texttt{ClassNotFoundException},
o JAR do Xposed esta no classpath --- o framework esta ativo.
\end{historybox}

\subsection{checkFrida() --- A Deteccao Mais Sofisticada}

\begin{historybox}
\textbf{Historia: Frida (Ole Andre V. Ravnas, Noruega, 2014)}

Ole Andre queria algo que o GDB (criado por Richard Stallman em 1986)
nao oferecia: injetar JavaScript num processo em execucao, sem pausar,
interativamente. Criou o Frida em 2014.

Arquitetura cliente-servidor: \texttt{frida-server} roda como root no device,
usa \texttt{ptrace()} --- chamada de sistema Unix criada em \textbf{1979}
no Unix V7 pelos engenheiros da Bell Labs --- para injetar \texttt{frida-agent.so}
no processo alvo via o dynamic linker.
\end{historybox}

\begin{lstlisting}
// Deteccao 1: probe de socket (150ms timeout)
Socket().use { socket ->
    socket.connect(InetSocketAddress("127.0.0.1", 27042), 150)
    signals += DetectionSignal(..., Severity.HIGH)
}

// Deteccao 2: /proc/net/tcp fallback
if (signals.none { it.description.contains("Frida server") }) {
    signals += checkFridaTcpTable("/proc/net/tcp")
    signals += checkFridaTcpTable("/proc/net/tcp6")
}
\end{lstlisting}

\begin{historybox}
\textbf{Bug encontrado e corrigido: porta Frida}

Durante o desenvolvimento foi encontrado um bug critico:

\begin{lstlisting}
// Codigo com bug:
val fridaPortHex = "699A"  // = 27034 decimal -- ERRADO

// Codigo correto:
val fridaPortHex = "69A2"  // = 27042 decimal -- CORRETO
\end{lstlisting}

\texttt{hex(27042) = 0x69A2}. O codigo anterior procurava a porta 27034
no \texttt{/proc/net/tcp}. O Frida escuta em 27042. O fallback inteiro
era ineficaz --- silenciosamente, sem erro.

Os testes passavam porque testavam a logica do parser com linhas
sinteticas, nao a conversao do numero da porta.
\end{historybox}

\subsubsection{O formato do /proc/net/tcp}

\begin{historybox}
\textbf{Historia: Big-Endian vs Little-Endian (1978 vs 1974)}

O arquivo \texttt{/proc/net/tcp} mistura duas convencoes historicas:

\begin{lstlisting}
local_address = "0100007F:69A2"
                 ^^^       ^^^
              little-endian  big-endian
              (Intel 1978)   (TCP/IP 1974)
\end{lstlisting}

O IP \texttt{0100007F} usa little-endian porque o kernel Linux roda em
processadores Intel x86, que desde 1978 armazenam bytes do menos significativo
primeiro. A porta \texttt{69A2} usa big-endian (network byte order) porque o
protocolo TCP foi definido por \textbf{Vint Cerf} e \textbf{Bob Kahn} em 1974
com bytes do mais significativo primeiro.

O estado \texttt{0A} = 10 decimal = \texttt{TCP\_LISTEN} --- a maquina de
estados TCP foi formalizada por Cerf e Kahn no mesmo paper de 1974.

\textbf{Danny Cohen} da USC cunhou os termos ``big-endian'' e ``little-endian''
em 1980 no paper \textit{``On Holy Wars and a Plea for Peace''}, referenciando
o livro \textit{As Viagens de Gulliver} de Jonathan Swift.
\end{historybox}

\newpage

% ─── ARQUIVO 8 ───────────────────────────────────────────────────────────────
\section{RootDetector.kt --- Detectando Root}

\begin{lstlisting}
class RootDetector(
    private val fileProbe: (String) -> Boolean = { File(it).exists() },
    private val writableProbe: (String) -> Boolean =
        { File(it).run { exists() && canWrite() } },
    private val packageChecker: (Context, String) -> Boolean = { ctx, pkg ->
        runCatching { ctx.packageManager.getPackageInfo(pkg, 0); true }
            .getOrDefault(false)
    },
) : EnvironmentAnalyzer
\end{lstlisting}

\begin{historybox}
\textbf{Historia: Root no Unix e Android (1969 -- hoje)}

Ken Thompson e Dennis Ritchie criaram o usuario \textbf{root} (UID 0) no Unix
em 1969 --- acesso irrestrito, sem limitacoes.

No Android, root nao existe por padrao. A Google removeu deliberadamente.
Mas a comunidade criou ferramentas:

\medskip
\textbf{SuperSU (Chainfire, Holanda, 2012)} --- o padrao-ouro por anos.
Mostrava um popup pedindo permissao cada vez que um app tentava elevar privilegios.
Em 2017, Chainfire vendeu para uma empresa chinesa (CCMT). A comunidade migrou.

\medskip
\textbf{Magisk (topjohnwu / Wang Guang, Taiwan, 2016)} --- systemless root.
O Android tem \textbf{dm-verity}: verificacao criptografica de cada bloco da
particao \texttt{/system} no boot. Root tradicional quebrava o dm-verity.
Magisk monta \texttt{tmpfs} sobre \texttt{/system} sem modificar o original.

\medskip
\textbf{KingRoot, Kingo Root (China, 2013-2014)} --- populares no mercado
asiatico por rootear sem desbloquear o bootloader, explorando vulnerabilidades
do kernel.
\end{historybox}

\subsection{Camada 1: checkSuBinaries()}

\begin{lstlisting}
internal fun checkSuBinaries(): List<DetectionSignal> =
    DetectionSignatures.SU_PATHS.filter { fileProbe(it) }.map { path ->
        DetectionSignal(SignalCategory.ROOT, "Root binary found: $path", HIGH)
    }
\end{lstlisting}

O binario \texttt{su} (switch user) e um dos mais antigos do Unix.
Em Android de producao, nao existe. Ferramentas de root o instalam em:
\texttt{/system/bin/su}, \texttt{/system/xbin/su}, \texttt{/sbin/su}, \texttt{/su/bin/su}.

\subsection{Camada 3: checkDangerousProps()}

\begin{lstlisting}
val explanation = when (prop) {
    "ro.debuggable" -> "ro.debuggable=1 habilita JDWP em qualquer processo..."
    "ro.secure"     -> "ro.secure=0 concede shell root via ADB sem autenticacao..."
    "ro.build.type" -> "Build type 'eng' ou 'userdebug' indica ambiente nao-producao..."
}
\end{lstlisting}

\begin{historybox}
\textbf{Historia: JDWP e ro.debuggable}

\textbf{JDWP} (Java Debug Wire Protocol) foi criado pela Sun Microsystems em 1998
como parte do Java Platform Debugger Architecture. Permite que qualquer processo
seja depurado via rede --- leitura de memoria, modificacao de variaveis em runtime.

\texttt{ro.debuggable=1} habilita JDWP em qualquer processo do sistema,
incluindo apps com \texttt{android:debuggable=false}. Em builds de producao sempre
e 0. Presente apenas em builds de engenharia ou dispositivos rootados.

\texttt{ro.build.type} tem tres valores:
\texttt{user} (producao), \texttt{userdebug} (depuracao habilitada), \texttt{eng} (root automatico).
\end{historybox}

\subsection{Camada 4: checkWritableSystemPaths()}

\begin{lstlisting}
internal fun checkWritableSystemPaths(): List<DetectionSignal> =
    DetectionSignatures.WRITABLE_SYSTEM_PATHS.filter { writableProbe(it) }.map { path ->
        DetectionSignal(SignalCategory.ROOT,
            "System path is writable -- indicates root mount: $path", HIGH)
    }
\end{lstlisting}

Em todo device de producao, \texttt{/system} e montado como \textbf{read-only}
pelo \texttt{init} no boot. O dm-verity garante isso criptograficamente.
\texttt{File("/system").canWrite() == true} e impossivel em device limpo.

\newpage

% ─── ARQUIVO 9 ───────────────────────────────────────────────────────────────
\section{SignatureAnalyzer.kt --- Verificando Integridade do APK}

\begin{lstlisting}
override fun analyze(context: Context): List<DetectionSignal> =
    checkCertificate(context) +
    checkInstallerSource(context) +
    checkPackageNameConsistency(context)
\end{lstlisting}

\begin{historybox}
\textbf{Historia: Criptografia de Chave Publica e X.509 (1976-1988)}

\textbf{Whitfield Diffie} e \textbf{Martin Hellman} publicaram em 1976
\textit{``New Directions in Cryptography''} --- inventando a criptografia de
chave publica. Voce tem duas chaves matematicamente ligadas: a privada (so voce
tem) assina, a publica (todos podem ter) verifica. Impossivel forjar sem a privada.

Em 1988, a \textbf{ITU-T} publicou o \textbf{X.509} --- formato padrao para
certificados digitais. Um certificado X.509 contem: subject (quem e), issuer
(quem emitiu), validade, chave publica e assinatura.

O Android usa X.509 para assinar APKs desde a versao 1.0.
\end{historybox}

\subsection{Camada 1: Debug Key Detection}

\begin{lstlisting}
val cert = CertificateFactory.getInstance("X.509")
    .generateCertificate(ByteArrayInputStream(sig.toByteArray())) as X509Certificate
val subject = cert.subjectDN.name
if (subject.contains("Android Debug", ignoreCase = true)) {
    signals += DetectionSignal(..., Severity.MEDIUM)
}
\end{lstlisting}

Quando voce instala o Android Studio, ele cria automaticamente um
\textbf{debug keystore} com subject sempre igual a \texttt{CN=Android Debug, O=Android, C=US}.
Todo desenvolvedor no mundo tem esse mesmo subject.

Quando um atacante usa \texttt{apktool} para modificar e reassinar um APK,
ele usa sua propria debug key. O subject \texttt{CN=Android Debug} no APK
de producao e o fingerprint mais obvio de um repackaging.

\subsection{Camada 2: SHA-256 Hash Comparison}

\begin{lstlisting}
private fun computeHash(signature: Signature, algorithm: String): String =
    MessageDigest.getInstance(algorithm)
        .digest(signature.toByteArray())
        .joinToString("") { "%02X".format(it) }
\end{lstlisting}

\begin{historybox}
\textbf{Historia: SHA-256 (NSA, 2001)}

SHA (Secure Hash Algorithm) foi desenvolvido pela \textbf{NSA} e publicado
pelo \textbf{NIST} em 1993. O SHA-256, parte da familia SHA-2, saiu em 2001.

Propriedade fundamental: qualquer mudanca no input, por menor que seja,
muda completamente o output. Trocar um bit no APK muda o hash inteiramente.

\texttt{MessageDigest} foi introduzido no Java 1.1 em 1997 pela Sun Microsystems
como parte do Java Cryptography Architecture.
\end{historybox}

\subsection{A API Dividida: Android P (2018)}

\begin{lstlisting}
if (Build.VERSION.SDK_INT >= Build.VERSION_CODES.P) {
    // Android 9+ (2018): APK Signature Scheme v3, suporta rotacao de chaves
    pm.getPackageInfo(packageName, GET_SIGNING_CERTIFICATES)
        .signingInfo?.apkContentsSigners
} else {
    @Suppress("DEPRECATION")
    pm.getPackageInfo(packageName, GET_SIGNATURES).signatures
}
\end{lstlisting}

O Android 9 introduziu o APK Signature Scheme v3 com suporte a rotacao de chaves.
A API antiga nao suportava multiplos certificados historicos e foi deprecada.

\newpage

% ─── ARQUIVO 10 ──────────────────────────────────────────────────────────────
\section{DetectionSignatures.kt --- O Catalogo Centralizado}

\begin{lstlisting}
internal object DetectionSignatures {
    val KNOWN_CLONE_PACKAGES = listOf(
        "com.lbe.parallel.intl",   // Parallel Space
        "com.lody.virtual",        // VirtualApp
        "me.weishu.exp",           // VirtualXposed
        // ... 21 outros
    )
    val SU_PATHS = listOf("/system/bin/su", "/system/xbin/su", ...)
    val ROOT_PACKAGES = listOf("com.topjohnwu.magisk", "eu.chainfire.supersu", ...)
    val HOOK_PATHS = listOf("/system/framework/XposedBridge.jar", ...)
    val LEGITIMATE_INSTALLERS = setOf("com.android.vending", ...)
    // ...
}
\end{lstlisting}

\begin{historybox}
\textbf{Historia: Single Responsibility Principle (Robert C. Martin, 2000)}

\textbf{Robert C. Martin} (Uncle Bob) formalizou o \textbf{Single Responsibility
Principle} em 2000: cada modulo deve ter apenas uma razao para mudar.

Antes desse arquivo, cada analisador mantinha suas proprias listas de
assinaturas. Quando um novo clone app aparecia, voce precisava atualizar
varios arquivos. \texttt{DetectionSignatures} centraliza: so um arquivo
muda quando a threat intelligence muda.

\texttt{internal object} garante que o catalogo e visivel em todo o modulo
\texttt{:app} mas invisivel fora dele --- os dados de seguranca nao vazam
para dependencias externas.
\end{historybox}

\newpage

% ─── ARQUIVO 11 ──────────────────────────────────────────────────────────────
\section{AndroidManifest.xml --- A Configuracao}

\begin{lstlisting}
<queries>
    <package android:name="com.lbe.parallel.intl" />
    <package android:name="com.lody.virtual" />
    <!-- ... 22 outros clone apps -->
</queries>
\end{lstlisting}

\begin{historybox}
\textbf{Historia: Package Visibility Filtering no Android 11 (2020)}

Antes do Android 11, qualquer app podia chamar
\texttt{PackageManager.getInstalledPackages()} e ver todos os apps instalados.

O Google, preocupado com privacidade, introduziu o \textbf{Package Visibility
Filtering} no Android 11 (API 30, 2020): apps sao invisiveis ao PackageManager
a menos que explicitamente declarados no \texttt{<queries>}.

Sem essa lista, no Android 11+, o \texttt{CloneAnalyzer} ficaria completamente
cego para apps clone --- \texttt{getPackageInfo("com.lbe.parallel.intl")}
lancaria \texttt{NameNotFoundException} como se o app nao existisse.

A alternativa seria \texttt{QUERY\_ALL\_PACKAGES}, mas o Google rejeita
apps na Play Store que pedem isso sem justificativa forte.
\end{historybox}

\newpage

% ─── CONCEITOS TRANSVERSAIS ───────────────────────────────────────────────────
\section{Conceitos Transversais}

\subsection{O Padrao de Injecao de Dependencias}

Todos os analisadores usam o mesmo padrao:

\begin{lstlisting}
class RootDetector(
    private val fileProbe: (String) -> Boolean = { File(it).exists() },
    // ...
)
\end{lstlisting}

Em producao: funcoes reais do sistema (nao passa nada, usa defaults).
Nos testes: lambdas que simulam qualquer cenario sem precisar de device.

Esse padrao foi formalizado por \textbf{Michael Feathers} em 2004 e e a base
para todos os testes unitarios do projeto (usando \textbf{MockK}).

\subsection{runCatching e Result<T>}

\begin{lstlisting}
runCatching { parseTcpTable(File(path).readLines(), path) }
    .getOrDefault(emptyList())
\end{lstlisting}

Em C (Dennis Ritchie, 1972), erros eram tratados por valor de retorno.
Em C++ (Bjarne Stroustrup, 1983) e Java (James Gosling, 1995), vieram
as excecoes com \texttt{try/catch}.

Linguagens funcionais dos anos 80 (Haskell, ML) propuseram tratar o erro
como um \textbf{valor} --- \texttt{Result<T>} do Kotlin encapsula o resultado
como sucesso ou falha sem interromper o fluxo.

\subsection{O Binder IPC do Android}

\begin{historybox}
\textbf{Historia: OpenBinder e o protocolo de comunicacao do Android}

Em 1996, a \textbf{Be Inc.} criou o \textbf{OpenBinder} para o BeOS ---
comunicacao entre processos via mensagens estruturadas.
A Be Inc. faliu em 2001. \textbf{Dianne Hackborn}, que trabalhava la,
levou a ideia para o Google em 2005 e reimplementou como \textbf{Binder IPC}.

Todo acesso ao sistema Android (PackageManager, LocationManager, etc.)
passa pelo Binder. O Android coloca dentro de cada processo um proxy local:
\texttt{IPackageManager.Stub.Proxy} --- o nome segue o padrao \textbf{CORBA (1991)}
e \textbf{Java RMI (1997)}: Stub (lado servidor, desempacota) e Proxy
(lado cliente, empacota).

O VirtualApp substituia esses proxies via Reflection para interceptar
todas as chamadas de sistema --- sem root, apenas com APIs publicas do Java.
\end{historybox}

\newpage

% ─── BUG ENCONTRADO ──────────────────────────────────────────────────────────
\section{Bug Encontrado e Corrigido}

Durante o desenvolvimento foi identificado um bug critico no
\texttt{VirtualizationAnalyzer.parseTcpTable()}:

\begin{lstlisting}
// ANTES (com bug):
val fridaPortHex = "699A"   // 0x699A = 27034 decimal -- PORTA ERRADA

// DEPOIS (corrigido):
val fridaPortHex = "69A2"   // 0x69A2 = 27042 decimal -- PORTA CORRETA
\end{lstlisting}

\subsection{Como o bug foi encontrado}

\begin{lstlisting}[language=Python]
>>> hex(27042)         # porta real do Frida
'0x69a2'
>>> int('699A', 16)    # o que estava no codigo
27034                  # ERRADO
>>> int('69A2', 16)    # o que deveria estar
27042                  # CORRETO
\end{lstlisting}

\subsection{Por que os testes nao pegaram}

Os testes do \texttt{parseTcpTable} testavam a \textbf{logica do parser}
com linhas sinteticas (rows com \texttt{"69A2"} hardcoded nos dados de teste),
nao a \textbf{conversao numerica} da constante \texttt{FRIDA\_PORT} para hex.

O fallback inteiro era ineficaz em producao: procurava a porta 27034
enquanto o Frida escutava em 27042. Silenciosamente, sem excecao, sem aviso.

\subsection{Impacto}

Com o bug, um atacante usando \texttt{frida-server} com a porta padrao teria
o probe de socket falhar (se bloqueado) E o fallback de \texttt{/proc/net/tcp}
retornar falso negativo. Deteccao via TCP completamente cega.

\newpage

% ─── RESUMO FINAL ────────────────────────────────────────────────────────────
\section{Resumo: Decadas de Historia em 500 Linhas de Kotlin}

\begin{center}
\begin{tabular}{p{3cm}p{4cm}p{6cm}}
\toprule
\textbf{Tecnica} & \textbf{Origem} & \textbf{Como usada no projeto} \\
\midrule
UIDs Unix & Bell Labs, 1969 & detecta clone apps pelo UID offset \\
\texttt{/proc} filesystem & Bell Labs, 1984 & le \texttt{maps} e \texttt{net/tcp} \\
ELF shared libs & Unix System V, 1988 & paths baked-in nos \texttt{.so} \\
X.509 certificates & ITU-T, 1988 & verifica assinatura do APK \\
Java Reflection & Sun, 1997 & acessa APIs internas do Android \\
SHA-256 & NSA/NIST, 2001 & hash do certificado de assinatura \\
Binder IPC & Be Inc./Google, 1996-2005 & base do VirtualApp que detectamos \\
\texttt{ptrace()} & Bell Labs, 1979 & base do Frida que detectamos \\
big/little-endian & Intel/TCP, 1974-1978 & parse do \texttt{/proc/net/tcp} \\
Android multi-user & Google, 2012 & UID / 100.000 detecta clone \\
Kotlin coroutines & JetBrains, 2018 & analisadores em paralelo \\
Bayesian probability & Thomas Bayes, 1763 & \texttt{emulatorProbability} \\
\bottomrule
\end{tabular}
\end{center}

\vspace{1cm}

\begin{center}
\textit{``A melhor deteccao de seguranca nao inventa novos truques.\\
Ela usa as limitacoes estruturais que os atacantes nao conseguem esconder.''}
\end{center}


\newpage

% ─── ROTEIRO DE APRESENTACAO ─────────────────────────────────────────────────
\section{Roteiro de Apresentacao}

\begin{mdframed}[backgroundcolor=codebg, linecolor=codeframe, linewidth=1pt,
  innerleftmargin=14pt, innerrightmargin=14pt, innertopmargin=10pt, innerbottommargin=10pt]
{\small\textit{Legenda: \textbf{[FALAR]} = o que dizer em voz alta.
\textbf{[MOSTRAR]} = arquivo ou funcao para abrir no IDE.
\textbf{[DURAÇÃO]} = tempo estimado de cada bloco.}}
\end{mdframed}

\vspace{0.5cm}

% ── ABERTURA ─────────────────────────────────────────────────────────────────
\subsection*{1. Abertura (1 minuto)}

\textbf{[FALAR]}

\begin{quote}
``Antes de mostrar o codigo, deixa eu contextualizar o problema que o projeto resolve.
Aplicacoes financeiras precisam saber se estao rodando num dispositivo legitimo ---
ou dentro de um ambiente controlado por um atacante.
Emuladores, apps de clone, ferramentas de hooking como Frida e Xposed, e APKs
reempacotados sao os vetores mais comuns de fraude mobile hoje.
\textit{Portanto}, o GuardApp e um SDK Android que detecta esses ambientes
em tempo de execucao, sem depender de servidor externo.''
\end{quote}

\vspace{0.3cm}

% ── ARQUITETURA ──────────────────────────────────────────────────────────────
\subsection*{2. Arquitetura Geral (2 minutos)}

\textbf{[MOSTRAR]} \texttt{EnvironmentGuard.kt}

\textbf{[FALAR]}

\begin{quote}
``\textit{Para comecar}, a decisao de design mais importante do projeto esta aqui,
nessa unica interface de 3 linhas ---
\texttt{EnvironmentAnalyzer} define o contrato que todo detector deve seguir.
\textit{A partir disso}, o \texttt{EnvironmentGuard} orquestra todos os analisadores
em paralelo usando coroutines Kotlin ---
cada um roda no \texttt{Dispatchers.Default} simultaneamente,
e qualquer falha individual e capturada silenciosamente com \texttt{runCatching}.
\textit{Dessa forma}, um analisador que crasha nao derruba o relatorio inteiro.
\textit{Alem disso}, adicionar um novo detector no futuro e apenas instanciar
uma nova classe na lista --- nenhuma outra classe precisa mudar.''
\end{quote}

\vspace{0.3cm}

% ── EMULADOR ─────────────────────────────────────────────────────────────────
\subsection*{3. EmulatorAnalyzer --- Detectando AVD e QEMU (2 minutos)}

\textbf{[MOSTRAR]} \texttt{EmulatorAnalyzer.kt} $\to$ \texttt{checkBuildProperties()} e \texttt{checkCpuInfo()}

\textbf{[FALAR]}

\begin{quote}
``\textit{Comecando pelo} detector de emulador, ele usa quatro camadas independentes.
\textit{Primeiro}, verifica as propriedades do \texttt{android.os.Build} ---
valores como \texttt{Build.HARDWARE = "goldfish"} ou \texttt{Build.FINGERPRINT}
comecando com \texttt{"generic"} so existem no emulador oficial do Android,
o AVD, que usa o QEMU criado por Fabrice Bellard em 2003.
\textit{Em seguida}, acessa propriedades internas do sistema via Reflection ---
como \texttt{ro.kernel.qemu=1} --- que nao estao no SDK publico mas
o kernel do emulador seta automaticamente.
\textit{Por fim}, tem um detalhe tecnico que vale mencionar:
a string \texttt{"goldfish"} nunca aparece como literal no APK compilado ---
ela e construida em runtime a partir de um \texttt{charArrayOf}.
\textit{Isso porque} qualquer um pode rodar \texttt{strings app.apk | grep goldfish}
e descobrir exatamente o que o app esta checando.
Construir em runtime impede extracao estatica.''
\end{quote}

\vspace{0.3cm}

% ── CLONE ────────────────────────────────────────────────────────────────────
\subsection*{4. CloneAnalyzer --- Detectando VirtualApp e Parallel Space (2 minutos)}

\textbf{[MOSTRAR]} \texttt{CloneAnalyzer.kt} $\to$ \texttt{checkSecondaryUserEnvironment()}

\textbf{[FALAR]}

\begin{quote}
``\textit{Passando para} o detector de clone apps, esse e o mais elegante do projeto.
\textit{Para entender}, preciso dar um contexto historico rapido:
em 2015, um estudante chines chamado Lody criou o VirtualApp ---
um framework que cria um Android dentro do Android interceptando chamadas Binder,
que e o protocolo IPC do Android criado pela Dianne Hackborn.
\textit{Consequentemente}, surgiu o Parallel Space e dezenas de similares,
usados em fraudes bancarias a partir de 2017.
\textit{A deteccao mais confiavel aqui} e o offset de UID:
o Android aloca UIDs em blocos de 100.000 por usuario desde 2012,
decisao de engenharia feita para tablets de familia.
\textit{Portanto}, qualquer app rodando em clone tera
\texttt{myUid() / 100.000 != 0} --- impossivel de falsificar sem root.
\textit{Alem disso}, o detector verifica o \texttt{/proc/self/maps}:
mesmo que o atacante renomeie o APK clone com apktool,
as bibliotecas \texttt{.so} compiladas tem os paths originais baked-in pelo linker ---
o kernel nao mente sobre o que esta carregado na memoria.''
\end{quote}

\vspace{0.3cm}

% ── VIRTUALIZATION ────────────────────────────────────────────────────────────
\subsection*{5. VirtualizationAnalyzer --- Detectando Frida e Xposed (3 minutos)}

\textbf{[MOSTRAR]} \texttt{VirtualizationAnalyzer.kt} $\to$ \texttt{checkXposedClassLoader()} e \texttt{parseTcpTable()}

\textbf{[FALAR]}

\begin{quote}
``\textit{Esse e o analisador mais complexo}, e tambem onde encontramos
e corrigimos um bug real durante o desenvolvimento.
\textit{Comecando pelo} Xposed: rovo89, desenvolvedor alemao, em 2012
substituiu o \texttt{app\_process} --- o binario que inicia todo app Android ---
por uma versao que injeta o \texttt{XposedBridge.jar} no classpath do Zygote.
\textit{Como resultado}, toda classe do Xposed fica acessivel em qualquer processo.
\textit{Portanto}, a deteccao e simples: se \texttt{Class.forName()} encontra
\texttt{"de.robv.android.xposed.XposedBridge"} sem lancar
\texttt{ClassNotFoundException}, o framework esta ativo.

\textit{Para o Frida}, criado pelo noruegues Ole Andre em 2014,
a deteccao usa duas estrategias.
\textit{Primeiro}, probe de socket na porta 27042 com timeout de 150ms.
\textit{Caso isso falhe}, um fallback que le \texttt{/proc/net/tcp}
procurando a porta em estado \texttt{TCP\_LISTEN}.

\textit{E aqui entra o bug que encontramos}:
o codigo original usava \texttt{"699A"}, que e 27034 decimal.
\textit{Porem}, a porta do Frida e 27042, que em hex e \texttt{0x69A2}.
\textit{Por consequencia}, o fallback inteiro era cego --- procurava
a porta errada silenciosamente, sem excecao, sem aviso nos logs.
Os testes passavam porque testavam a logica do parser
com linhas sinteticas, nao a conversao do numero em si.''
\end{quote}

\vspace{0.3cm}

% ── ROOT ─────────────────────────────────────────────────────────────────────
\subsection*{6. RootDetector --- Detectando Root e Magisk (1 minuto)}

\textbf{[MOSTRAR]} \texttt{RootDetector.kt} $\to$ construtor com lambdas injetaveis

\textbf{[FALAR]}

\begin{quote}
``\textit{O detector de root usa quatro estrategias} complementares:
binarios \texttt{su} em paths conhecidos, apps de gerenciamento como Magisk,
propriedades de sistema como \texttt{ro.debuggable=1} e \texttt{ro.secure=0},
e verificacao se \texttt{/system} esta montado com escrita.
\textit{O detalhe arquitetural importante aqui} e o construtor com lambdas injetaveis.
\textit{Isso permite} testar cada verificacao sem precisar de device fisico rootado ---
o teste passa uma lambda que simula o sistema de arquivos e o
\texttt{PackageManager} retornando os valores desejados.
\textit{Esse padrao} foi formalizado por Michael Feathers em 2004
e e a base de todos os testes unitarios do projeto.''
\end{quote}

\vspace{0.3cm}

% ── SIGNATURE ────────────────────────────────────────────────────────────────
\subsection*{7. SignatureAnalyzer --- Verificando Integridade do APK (1 minuto)}

\textbf{[MOSTRAR]} \texttt{SignatureAnalyzer.kt} $\to$ \texttt{checkCertificate()}

\textbf{[FALAR]}

\begin{quote}
``\textit{Por fim}, o detector de integridade do APK.
\textit{Quando um atacante} pega o APK, modifica com \texttt{apktool}
e reassina com a debug key propria, o certificado X.509 vai ter
\texttt{CN=Android Debug} no subject.
\textit{Alem disso}, o SHA-256 do certificado vai ser diferente do esperado.
\textit{Para producao}, basta configurar \texttt{EXPECTED\_CERT\_SHA256}
com o hash do certificado de release --- qualquer APK reempacotado
vai ter hash diferente e sera detectado com \texttt{Severity.HIGH}.
\textit{Complementarmente}, o analisador verifica o instalador do app:
se nao veio da Play Store ou de uma loja reconhecida,
e um sinal de distribuicao fora dos canais oficiais.''
\end{quote}

\vspace{0.3cm}

% ── ENCERRAMENTO ─────────────────────────────────────────────────────────────
\subsection*{8. Encerramento (1 minuto)}

\textbf{[FALAR]}

\begin{quote}
``\textit{Para resumir}, o projeto detecta cinco vetores de ataque ---
emulador, clone app, hooking, root e reempacotamento ---
usando uma arquitetura extensivel onde cada detector e independente
e roda em paralelo via coroutines.
\textit{Alem do codigo em si}, encontramos e corrigimos um bug real
de producao durante o desenvolvimento: o port hex do Frida estava errado
e o fallback era completamente ineficaz.
\textit{Do ponto de vista tecnico}, cada decisao de design tem uma justificativa ---
desde o \texttt{charArrayOf} para evasao de analise estatica
ate o padrao de injecao de dependencia para testabilidade sem device.
\textit{Estou a disposicao para perguntas sobre qualquer parte do projeto.}''
\end{quote}

\vspace{0.5cm}

\begin{mdframed}[backgroundcolor=histbg, linecolor=histframe, linewidth=1.5pt,
  innerleftmargin=14pt, innerrightmargin=14pt, innertopmargin=10pt, innerbottommargin=10pt]
\textbf{Tempo total estimado: 13 minutos.}
Deixe 5-7 minutos para perguntas.
Pontos que costumam gerar mais interesse: o bug do port hex,
o charArrayOf, e o UID offset de 100.000.
Se pedirem para mostrar um teste, abra \texttt{VirtualizationAnalyzerTest.kt}
ou \texttt{RootDetectorTest.kt} --- os mais didaticos.
\end{mdframed}

\end{document}
